% vim: ft=tex
% SECTION: DESIGN

I used Kicad\footnote{http://kicad.org} for schematic capture and board design.
All of the design files can be found on the LED controller page of the ACRIS
Project.\footnote{http://jwcxz.com/projects/acris/ledctrlr.php}

\subsection{LED Drivers}
    TI's TLC5940 is a device that reads serial data from some source and
    uses it to individually set the PWM duty cycle for 16 constant current
    sinks.  Each of the 16 output channels can sink up to a maximum of
    120mA.  Channels can be tied together in order to sink more current.

    There are several operational quirks regarding the TLC5940.  The
    chip requires an external clock source to set the PWM frequency; it
    does not have an auto-reloading timer, so the \pin{blank} pin must be
    twiddled every 4096 ticks of the \pin{gsclk} clock.  These chips are
    also extremely sensitive to static discharge and overvoltage on any
    pin, even the sink pins.  Finally, the \pin{iref} pin is connected to a
    resistor that sets the maximum current sink of the chip.  It is
    determined by equation~\ref{eqn:design:elec:iref}.

    \begin{align}
        \label{eqn:design:elec:iref}
        I_\textrm{max} = \frac{1.24}{R_\textrm{IREF}} \cdot 31.5
    \end{align}

    So, for RGB LEDs, we chose to use a $330\Omega$ resistor in order to
    yield a per-channel current of 118.3mA; we then tied three of these
    outputs together to yield around 355mA current sink capability.  For
    the UV and white LEDs, we tied 7 output channels together with
    $360\Omega$ resistors on the drivers to yield a per-channel current of
    108.5mA and a total current of 759.5mA.
    Section~\ref{sec:design:elec:leds} discusses the properties of the LEDs
    we chose.

    Sending data to the TLC5940 is relatively easy.  Each channel has
    12-bit resolution, so a total of $12\cdot16=192$ bits are sent via the
    serial input to control a single chip.  Furthermore, we chose to daisy
    chain the chips; that is, the serial output of the first chip feeds the
    serial input of the second chip, whose output then feeds the serial
    input of the third chip.  The \pin{xlat} pin is twiddled after
    $192\cdot3=576$ bits have been sent in.  This latches the data, causing
    the TLC5940 to drive the LED outputs with a different duty cycle.
    UV/strobes only use two TLC5940s, so a total of $192\cdot2=384$ bits
    are sent through these drivers.

    As noted in Section~\ref{sec:design:elec:ledctrl:comm}, the
    communications system only allows for 8-bit resolution in setting the
    PWM duty cycle of the LEDs.  This value is upconverted in order to fill
    the 12-bit range of the LED drivers.

\subsection{Microcontroller}
    The fundamental job of the microcontroller is to read incoming serial
    data, determine what action it needs to take based on that serial data,
    and, if necessary, update the LED output values on the TLC5940 chips.
    We chose the ATMEGA168P for its low cost and excellent functionality.
    Furthermore, we chose the version of the chip that was capable of using
    a 20MHz clock and drove it at this frequency.

    In addition to the code to handle the above behavior, the
    microcontroller handles timers and interrupts.  First, it
    generates a clock line for the TLC5940s' \pin{gsclk} signal at a rate
    of 1MHz.  Second, it generates a blank at the rate of
    $\frac{1\textrm{MHz}}{2^{12}}=244\textrm{Hz}$.  Thus, the PWM frequency
    of the TLC5940 is 244Hz.

    The microcontroller also uses an interrupt-driven UART framework.
    Received UART data is stored to a ring buffer of size 128 bytes.  The
    main code processes it as needed.

    Lastly, the microcontroller has a timeout timer; if no data is received
    after a few seconds, the timeout timer triggers and the microcontroller
    displays a simple moving pattern on the RGB displays.  However, due to
    a bug in the code in combination with the fixes we had to implement,
    this feature does not work correctly on almost all of the front panels
    and one of the rear panels.

    The main loop of the microcontroller begins by setting up all of the
    necessary timers and interrupts for the system.  It then displays a
    startup pattern, which varies between RGB and UV/strobe panels.  The
    RGB panels initially display four red pixels with the middle two being
    brighter than the outer two.  Then, this is cleared and the first pixel
    is set to be green.  Then, the address is displayed in binary format on
    the pixels in blue (the first pixel retains its green coloration and
    turns cyan when the high bit is 1).

    The UV/strobe displays its address by its three columns.  E.g., if the
    instrument address is \ttt{3}, then the middle column (the UV LEDs) and
    the rightmost column (the white LEDs on the right) will both be lit.
    In the event that more UV/strobes are added, addresses
    \ttt{8}--\ttt{15} will look brighter.

    Then, the system enables interrupts in order to start receiving data.
    When the buffer is not empty, the main loop consumes the characters in
    the buffer as it runs through a state machine that determines which
    action to take.  When an action's arguments have been fulfilled, the
    system drives the LEDs appropriately and begins processing the next
    received data.  The packet processing state machine is described by
    Figure~\ref{fig:design:elec:statemachine}.  Note that any instance of
    \ttt{0xAA} immediately returns the state back to its sync'd state,
    \ttt{SYNC}.  Therefore, LEDs cannot have a value of \ttt{170}; we
    decided that this was a good enough assumption to make so that our
    state machine never falls out of sync; the difference between \ttt{169}
    and \ttt{171} is not noticeable enough to need a \ttt{170}.  The only
    complication is that all software that controls the lighting system
    must take this into account when it prepares packets.

    \begin{rfigure*}
        \label{fig:design:elec:statemachine}
        \includegraphics[scale=1.0]{fig/photo.pdf}
        %\includegraphics[scale=0.6]{fig/fig-design-elec-statemachine.jpg}
        \caption{Packet-processing state machine}
    \end{rfigure*}

    The final result of processing a packet is to ultimately drive the LEDs
    with new levels.  So, when the state machine has finished processing a
    packet, it calls a wrapper function called \fn{led\_action()}, which
    then calls the appropriate action function to set the LEDs according to
    the arguments given.

    As noted earlier, the TLC5940s use 12-bit values to set the PWM duty
    cycle but the communications system only uses 8-bit color arguments.
    So, right now, we upconvert by simply shifting the 8-bit value left 4
    times.  However, we then need to store that data in an efficient
    manner, so we perform 12-bit packing in 8-bit arrays of values.  Each
    TLC5940 requires $16\cdot12=192\textrm{bits}$ of data to send, so by
    packing the 12-bit values into an array of 24 8-bit values, it will be
    easy to send all of this data directly to the LED drivers.

    A 12-bit number can be split into a byte and a nibble.  A left-tailed
    split is a split where the uppermost nibble of the number is the nibble
    and the two lower nibbles make up the byte (a right-tailed split uses
    the lowest nibble as the nibble).  So, given an array of 24 8-bit
    values $X$, we pack the 12-bit number associated with the TLC5940
    output channel $i$ (for $0\leq i \leq 15$) according to the following algorithm.

    First, define the base index $i_b$.  This is the index of $X$ that will store
    the byte of the 12-bit number.  It is dependent on whether $i$ is even
    or odd.

    \begin{equation}
        i_b = \left\{
            \begin{array}{ll}
                23 - \frac{3\cdot i}{2} & \textrm{$i$ even} \\
                23 - \frac{3\cdot i}{2} - 1 & \textrm{$i$ odd}
            \end{array}
            \right.
    \end{equation}
    
    Notice that 23 is the maximum index of $X$ and $\frac{3}{2}$ is the
    ratio $\frac{24}{16}$, which effectively scales $i$ to the new array
    length.  We pack in reverse order (i.e. $i=0 \Rightarrow i_b=23$) in
    order to make sending data to the TLC5940 happen in proper order.

    To complete the packing algorithm we do the following.  For even
    channels, we should store the two lower nibbles to the base index,
    $X[i_b]$ and the high nibble to $X[i_b-1]$'s bottom nibble.  For odd
    channels, we store the two higher nibbles to the base index and the low
    nibble to $X[i_b+1]$'s top nibble.

    Thus, for even values of $i$,
    
    \begin{align}
        X[i_b-1] &= val_i[11:8] \\
        X[i_b] &= val_i[7:0]
    \end{align}

    and for odd values of $i$,

    \begin{align}
        X[i_b] &= val_i[11:4] \\
        X[i_b+1][7:4] &= val_i[3:0]
    \end{align}

    For RGB instruments, we keep three of these arrays, one for each
    TLC5940.  For UV/strobes, we keep two, since UV/strobe boards only have
    two TLC5940s.  When the main loop calls a ``drive" function, the
    microcontroller sends data to the TLC5940s using its SPI peripheral.
    The \pin{sck} line from the microcontroller drives the \pin{sclk} line
    of the LED drivers, and the \pin{mosi} line feeds them the data.  The
    SPI controller can only send 8-bits at a time, but thanks to our
    packing scheme, doing so is just a matter of sending the 24 values in
    each array.  Since we tie the TLC5940s together serially, we just shift
    in all of the data for all of the drivers on the instrument and then
    twiddle the \pin{xlat} line on all of the devices to lock the values in
    and update the drivers.

\subsection{LED Controller Board}

    The majority of the LED controller's electronics are on a two-layer,
    4"$\times$2.5" board manufactured by Advanced
    Circuits.\footnote{http://4pcb.com}  Unfortunately, the board is not
    without its problems; we did some serious rework of the power system
    after we had already sent the boards out; this resulted in several
    bugs.  However, we were able to get around them all, albeit with some
    unpleasant wiring.  A full discussion of the problems we encountered
    with the electronics is in Section~\ref{sec:analysis:electronics}

    One of the primary features of the board is that it uses all
    surface-mount components in order to facilitate easy soldering.  We
    aimed to make \nhpls{} a learning tool; we had a very successful set of
    ``soldering parties" whereby students learned how to populate a board
    and solder well.

    The schematics, board design, and bill of materials we used is located
    at \textbf{TBD}.
    % TODO find a spot for schematics, board design, bom

    The relevant schematics and board layout diagram are also provided
    here, in Figures~\ref{fig:design:elec:schem} and
    \ref{fig:design:elec:brdlayout} respectively.

    \begin{rfigure*}
        \label{fig:design:elec:schem}
        \includegraphics[scale=1.0]{fig/photo.pdf}
        %\includegraphics[scale=0.6]{fig/fig-design-elec-schem.pdf}
        % TODO schematics
        \caption{LED controller board schematics}
    \end{rfigure*}

    \begin{rfigure*}
        \label{fig:design:elec:brdlayout}
        \includegraphics[scale=1.0]{fig/photo.pdf}
        %\includegraphics[scale=0.6]{fig/fig-design-elec-brdlayout1.pdf}
        %\includegraphics[scale=0.6]{fig/fig-design-elec-brdlayout2.pdf}
        % TODO brdlayout diagrams
        \caption{Layout of the LED controller board (top: RGB panel,
            bottom: UV/strobe panel)}
    \end{rfigure*}

    Regarding our schematics, we chose to make logic power entirely
    separate from LED power and decided to use point-to-point connectors to
    connect the positive LED rail to the LEDs.          

    In the top left-hand corner of the board is a spot for a barrel jack
    power connector.  We reconfigured the board to use LM7805s instead, as
    shown in Figure~\ref{fig:design:elec:barrelrework}.  The LEDs connect
    to the output ports around the LED drivers on the right-hand side of
    the board.

    \begin{rfigure}
        \label{fig:design:elec:barrelrework}
        \includegraphics[scale=1.0]{fig/photo.pdf}
        %\includegraphics[scale=0.6]{fig/fig-design-elec-barrelrework.pdf}
        % TODO barrel jack rework
        \caption{Barrel jack $\rightarrow$ LM7805 rework}
    \end{rfigure}


